% 2012 山东卷理 第9题 选项C：函数图象（错误选项，按原图描摹的示意图）
% 原图零点与极值不等距、外侧振幅增大；此错误选项无可反推公式，保留实际描摹路径。
\usetikzlibrary{arrows.meta}
\begin{tikzpicture}[x=1.7cm,y=0.71cm,line width=.6pt,font=\normalsize,>=Stealth]
\def\cCright{(0.09,2.70) (0.14,2.40) (0.19,1.30) (0.24,0.55) (0.29,0.00) (0.34,-0.85) (0.39,-1.25) (0.44,-1.30) (0.49,-1.15) (0.54,-0.75) (0.59,-0.25) (0.64,0.30) (0.69,0.62) (0.74,0.80) (0.79,0.72) (0.84,0.45) (0.89,0.20) (0.94,0.00) (0.99,-0.33) (1.04,-0.20) (1.09,0.10) (1.14,0.55) (1.19,1.00) (1.24,1.55) (1.29,2.15) (1.34,2.58) (1.35,2.70)}
\def\cCleft{(-1.35,-2.70) (-1.34,-2.58) (-1.29,-2.15) (-1.24,-1.55) (-1.19,-1.00) (-1.14,-0.55) (-1.09,-0.10) (-1.04,0.20) (-0.99,0.33) (-0.94,0.00) (-0.89,-0.20) (-0.84,-0.45) (-0.79,-0.72) (-0.74,-0.80) (-0.69,-0.62) (-0.64,-0.30) (-0.59,0.25) (-0.54,0.75) (-0.49,1.15) (-0.44,1.30) (-0.39,1.25) (-0.34,0.85) (-0.29,0.00) (-0.24,-0.55) (-0.19,-1.30) (-0.14,-2.40) (-0.09,-2.70)}
\coordinate (O) at (0,0);
\coordinate (Xend) at (2.5,0);
\coordinate (Yend) at (0,4.5);
\begin{scope}
  \clip (-2.2,-4.2) rectangle (2.5,4.2);
  \draw plot[smooth] coordinates {\cCleft};
  \draw plot[smooth] coordinates {\cCright};
\end{scope}
\draw[->] (-2.2,0)--(Xend) node[right,inner sep=0pt,xshift=2pt] {$x$};
\draw[->] (0,-4.5)--(Yend) node[above,inner sep=0pt,yshift=2pt] {$y$};
\node[below left,inner sep=1pt] at (O) {$O$};
\end{tikzpicture}
